% !TeX root = ../papers/nc-p3.tex
\section{四母線の重複と長方形配置の縮約点での余核}
\label{sec:sectional-hyperelliptic-grid}
\begin{theorem}\label{thm:sectional-hyperelliptic-grid}
定理\ref{thm:sectional-hyperelliptic-norm}の十八次の逆系 $A=2L+2F$ は
滑らかな二次曲面でも二次錐の解消でも生じない。
\end{theorem}
\begin{proof}[四母線上の局所表示]
前定理から $Z=(a,h)$、$h$ は底の二項四次式である。記号 $W,\Gamma,Y,\xi$ を引き続き使う。
$\Gamma\to Y\subset\mathbb P(W)$ は、高々四つの頂点を除いて同型である。
実際 $h$ の根の近くでは、相対二次式の基底を $a,b,c$ に取れ、
$W$ の基底は $a,hb,hc$ である。相対二次埋込みの座標 $[a:b:c]$ から
$[z_0:z_1:z_2]=[a:hb:hc]$ への射のグラフは $(a,h)$ のブローアップである
（基点では $b,c$ の一方が単元）。$z_1,z_2$ の一方が非零なら
$[h z_0:z_1:z_2]$ で逆座標が復元され、頂点 $[1:0:0]$ の外で同型となる。
これは $h$ の重根でも成立する。例外入力点のうち垂直なものは従って高々四個である。

$V\otimes\mathcal O\to W$ の像束の次数を $k$ とする。
一般には階数三、各繊維でも階数二以上である。階数一以下なら入力線形射影の中心が
相対平面内の直線を含み、必ず $Y$ の二次曲線に会ってしまうからである。
階数が落ちる底点は $4-k$ 個以下で、その平面の入力像は一本の直線。
飽和前の像束それ自身は $V$ の商束なので、その射影平面束
$\Sigma'\subset\mathbb P^1\times\mathbb P^3$ は二次数 $(k,1)$ の方程式 $g_x(t)$ で定まる。
$k=0$ は全入力像が一平面となり不可能。$\Gamma$ のそこへの像は類 $(8-2k,2)$ の整因子 $Y'$ である。
例外の入力直線の外では底の各該当平面が同型に写り、
さらに垂直縮約の像を除けば上の局所表示から $\Gamma$ と $Y'$ が一致する。
水平縮約は $k<4$ なら各例外入力直線に属する。

$k\le2$ なら例外直線、四垂直点、$g_x$ が恒等零となる共通部分を除いて入力繊維長は高々二。
共通部分は高々直線、例外直線は高々三本なので、これらには $3\cdot3+4+3<20$ 印しか載らない。
$k=3$ なら例外直線と垂直点の外に少なくとも $20-3-4=13$ 印がある。
ここで $M_D=f_*L$ により各繊維長は三である。
前定理の $(2,2)$ 式の三係数は全てこれらの印で消える。$g_x$ が恒等零なら
長さ三は二項二次式だけでは切れないので、この場合を別に加える必要はない。
互いに素な二係数の線形結合に対する上界十二に反する。

$k=4$ なら $V$ は $W$ を全ての繊維で生成し、$\Sigma'=\mathbb P(W)$、$Y'=Y$。
$g_x$ が恒等零となる入力線形空間上では $Y$ の式は底によらない定数係数の二次式である。
その零点は水平縮約の入力点で、高々二個である（共通部分が直線で制限が恒等零なら
その部分スクロールが整因子 $Y$ 自身となってしまう）。
その他の有限な印は前定理の四点因子の跡による同じ二次式
$\det l(x)=0$ 上にある。ここで $l(x)$ は二項二次式上の汎関数を
二行対称行列として書いたものである。引戻し二次式が恒等零なら共通の固定根の
入力平面に十四印以上が入って不可能。二次式には高々十四印なので、
四つの異なる垂直縮約点と二つの異なる水平縮約点が全て印である場合だけが残る。
\end{proof}

\begin{proof}[基点と水平切断]
四本の母線との交長が各二なので $h$ は単根のみを持つ。
前定理の通り $W=\mathcal O^2\oplus\mathcal O(4)$。各水平縮約上では一般の母線から
一意の点が写るので、$S_0$ 上の像は底の切断である。
この切断上の基点の総長は高々四である。実際 $Z\subset(h=0)$ で、
滑らかな切断上の $h$ の零因子の次数が四である。
従って $A$ による切断の次数も高々四でなければ入力で潰れない。
$\mathbb F_2$ で負切断以外の切断は $E+bF$, $b\ge2$ で $A(E+bF)=2+2b\ge6$。
水平縮約を二つ持てない。$\mathbb F_0$ では切断の型 $(1,b)$ の次数は $4+2b$ なので、
両方とも母線パラメータと独立な一定の点であり、四基点繊維のそれぞれに載る。
ゆえに $Z$ はこの二点と四つの底点の直積、全て単純である。

一定二点の座標を $s_0=0,s_1=0$ と選ぶ。$I_Z A$ の完全切断は
$u=h s_0^2$, $v=h s_1^2$ および $z S_4$（$z=s_0s_1$、$S_j$ は底の二項 $j$ 次式）で張られる。
実際両軸の係数四次式は四根で消えるからである。入力座標は
\[
      x_0=u+z b,\quad x_1=v+z c,\quad x_2=z r,\quad x_3=z r',
      \qquad b,c,r,r'\in S_4 .
\]
$r,r'$ は $k=4$ により共通根を持たない。二つの水平入力点は $e_0,e_1$ で、
$Y$ は $q=(x_0-zb)(x_1-zc)-z^2h^2=0$ の因子である。
\end{proof}

\begin{proof}[縮約点での二行表示]
$\Sigma=\mathbb P(W)$ から入力への射を $p$、$\Gamma\to Y$ を $g$ とする。
$g_*L$ は、次の行列の余核を底から $\mathcal O(1)$ で捻った層である。
\[
  T=\begin{pmatrix}x_1-zc&-zh\\-zh&x_0-zb\end{pmatrix}:
       \mathcal O_\Sigma(-\xi)^2\longrightarrow\mathcal O_\Sigma^2 .
\]
実際二生成元は母線座標 $s_0,s_1$ に評価でき、関係は恒等的に成り立つ。
$T$ の階数零は $u=v=0$, $h=0$ の有限個の頂点だけである。
その外では余核は $Y$ 上の線束で $s_0,s_1$ の評価は同型である。
余核は深さ二で $g_*L$ はtorsion-freeゆえ有限集合を越えても同型となる。
上は捻る前には $(1,0)$ の母線座標束で、底から捻れば核束 $L=(1,1)$ を得る。

$M_D=p_*g_*L$ の $e_0$ での繊維を計算する。
$\Sigma$ は $(4,1)$ 入射なので、上の二項を底から捻って得る
$\mathcal O_\Sigma(F)$ とその入力捻りは $p$ の第一高次直像が零である
（積空間上の入射の完全列で、底の $\mathcal O(1)$ は一次コホモロジーを持たない）。
よって押出し行列の余核が $M_D$ である。
$x_0=1$ のチャートで $\Sigma$ は $x_1$ のアフィン直線と、
底上の $\mathcal O(-4)$ の全空間の積である。後者の繊維座標を $z$ とすれば
入力の法座標は $z r,z r'$。従って $p_*\mathcal O_\Sigma(F)$ の $e_0$ での繊維は
\[
      P=S_1\ \oplus\ z\,S_5/(r,r')_5 .
\]
これは全空間の切断 $\bigoplus_{j\ge0} z^j S_{1+4j}$ を $z r,z r'$ で割ったもので、
二つの互いに素な四次式の商のsocle次数六より $j\ge2$ は消える。
両直和因子の次元は二である。
任意の $z b$ の作用は第一部分から第二部分への $b$ 倍だけで、これらの作用の積は零である。
$T(e_0)$ の単元 $1-zb$ を消去すると、残る作用は $-zc$。
従って $M_D(e_0)$ の次元は $4-\operatorname{rank}(c:S_1\to S_5/(r,r')_5)$ である。

この二行行列の階数は、前定理の境界汎関数 $l(e_0)$ の二項二次形式としての階数に一致する。
実際 $bc-h^2\in(r,r')_8$ だから、$q$ の $z^2$ の項は $x_2 z b_2+x_3 z b_3$ と書け、
境界で $e_0$ に残るのは $-zc$ の項だけである。
$z c$ の境界は $r,r'$ の鉛筆の列における $c$ の境界
$\partial c\in H^1(\mathcal O(-4))$ と符号を除いて一致する。
二チャート $r\ne0,r'\ne0$ で $z c$ を $x_2 c/r,x_3 c/r'$ として持ち上げれば、
差を入射式 $x_2r'-x_3r$ で割ったコサイクルは $c/(rr')$ である。
一方 $S_5/(r,r')_5\simeq H^1(\mathcal O(-3))$ で、$c$ 倍の写像の対は
$(s,s')\mapsto\langle\partial c,ss'\rangle$ である。従って階数の一致を得る。
$e_0$ が印なら $\dim M_D(e_0)=3$、よって $\det l(e_0)=0$。
$e_1$ でも $b,c$ を交換して同じである。
結局、高々四つの垂直点以外の十六印以上が同じ非零二次式に入る。上界十四に反する。
\end{proof}
\path{computations/verify_sectional_hyperelliptic_grid.py} は局所Artin環での二行表示、
境界の二項形式および両核曲面の縮約数を検算する。
